\documentclass[11pt]{article}

\usepackage{amsmath,amssymb,amsthm,mathtools}
\usepackage[margin=1in]{geometry}

\newtheorem{theorem}{Theorem}
\newtheorem{lemma}[theorem]{Lemma}
\newtheorem{corollary}[theorem]{Corollary}
\newtheorem{remark}[theorem]{Remark}
\newtheorem{definition}[theorem]{Definition}

\begin{document}

\section{Closure under full whiskering}

For a finite graph $H$ and $0\le p<1$, define
\[
    \Phi_H(p)
    :=
    (1-p)^{v(H)}
    \chi_H\!\left(\frac{1}{1-p}\right).
\]
Since $\chi_H$ is monic of degree $v(H)$, the right-hand side is
a polynomial in $1-p$, and therefore extends continuously to $p=1$.
We use this continuous extension when $p=1$.

Throughout, a graphon means a symmetric, jointly measurable function
$W:\Omega^2\to[0,1]$ on an arbitrary probability space
$(\Omega,\mathcal A,\mu)$.  Write
\[
    p=t(K_2,W)
      =
      \int_{\Omega^2}W(x,y)\,d\mu(x)d\mu(y)
\]
and
\[
    d_W(x):=\int_\Omega W(x,y)\,d\mu(y)
\]
for its degree function.

\begin{definition}[Full whiskering]
Let $F$ be a finite graph with vertex set $V(F)$.
The \emph{full whiskering} of $F$, denoted by $\operatorname{Wh}(F)$,
is obtained by adjoining, for each $v\in V(F)$, a new vertex $v'$
and the single new edge $vv'$.
Thus the newly added vertices are pairwise distinct and each of them
has degree one.

If $n=v(F)$, then
\[
    v(\operatorname{Wh}(F))=2n.
\]
\end{definition}

We first record the elementary path inequality needed in the proof.

\begin{lemma}\label{lem:P4}
Let $W$ be a graphon of edge density $p$. Let $P_4$ denote the path
on four vertices, hence with three edges. Then
\[
    t(P_4,W)\ge p^3.
\]
Equivalently,
\[
    \int_{\Omega^2}
        W(x,y)d_W(x)d_W(y)\,d\mu(x)d\mu(y)
    \ge p^3.
\]
\end{lemma}

\begin{proof}
Write $d=d_W$. Since $W\ge0$, if $d(x)=0$, then
\[
    W(x,y)=0
\]
for almost every $y$.  Tonelli and symmetry therefore show that the
set of pairs $(x,y)$ for which $W(x,y)>0$ but $d(x)d(y)=0$ has
$(\mu\otimes\mu)$-measure zero.  Hence, with the convention that
quotients by $d(x)=0$ are set equal to zero, we have almost everywhere
\[
\begin{aligned}
    W(x,y)
    &=
    \bigl(W(x,y)d(x)d(y)\bigr)^{1/3}
    \left(\frac{W(x,y)}{d(x)}\right)^{1/3}
    \left(\frac{W(x,y)}{d(y)}\right)^{1/3}.
\end{aligned}
\]
H\"older's inequality therefore gives
\[
\begin{aligned}
    p^3
    &=
    \left(\int_{\Omega^2}W(x,y)\,d\mu(x)d\mu(y)\right)^3
    \\
    &\le
    \left(
      \int_{\Omega^2}
      W(x,y)d(x)d(y)\,d\mu(x)d\mu(y)
    \right)
    \left(
      \int_{\Omega^2}
      \frac{W(x,y)}{d(x)}\,d\mu(x)d\mu(y)
    \right)
    \left(
      \int_{\Omega^2}
      \frac{W(x,y)}{d(y)}\,d\mu(x)d\mu(y)
    \right).
\end{aligned}
\]
For the second factor,
\[
\begin{aligned}
    \int_{\Omega^2}
      \frac{W(x,y)}{d(x)}\,d\mu(x)d\mu(y)
    &=
    \int_{\{x:d(x)>0\}}
       \frac{1}{d(x)}
       \left(\int_\Omega W(x,y)\,d\mu(y)\right)\,d\mu(x)
    \\
    &=
    \mu\{x:d(x)>0\}
    \le 1.
\end{aligned}
\]
The third factor is bounded by $1$ in the same way. Hence
\[
    p^3
    \le
    \int_{\Omega^2}
      W(x,y)d(x)d(y)\,d\mu(x)d\mu(y).
\]
Finally, integrating the two endpoints of a three-edge path first gives
\[
    t(P_4,W)
    =
    \int_{\Omega^2}
      W(x,y)d(x)d(y)\,d\mu(x)d\mu(y),
\]
which proves the claim.
\end{proof}

The key observation is that full whiskering naturally changes the
underlying vertex measure to the degree-biased measure.

\begin{theorem}[Full-whiskering closure]\label{thm:whiskering}
Let $F$ be a finite graph with
\[
    n=v(F),
\]
and let $p_0>0$. Assume the following two properties.

\begin{enumerate}
    \item For every graphon $V$ on every probability space whose edge
    density $s$ satisfies
    \[
        s\in[p_0,1],
    \]
    one has
    \[
        t(F,V)\ge \Phi_F(s).
    \]

    \item The function
    \[
        s\longmapsto \Phi_F(s)
    \]
    is nondecreasing on $[p_0,1]$.
\end{enumerate}

Then, for every graphon $W$ of edge density
\[
    p=t(K_2,W)\ge p_0,
\]
one has
\[
    t(\operatorname{Wh}(F),W)
    \ge
    \Phi_{\operatorname{Wh}(F)}(p).
\]
More precisely,
\[
    t(\operatorname{Wh}(F),W)
    \ge
    p^n\Phi_F(p),
\]
and
\[
    \Phi_{\operatorname{Wh}(F)}(p)
    =
    p^n\Phi_F(p).
\]
\end{theorem}

\begin{proof}
Let
\[
    d(x):=d_W(x)=\int_\Omega W(x,y)\,d\mu(y).
\]
Since
\[
    \int_\Omega d(x)\,d\mu(x)=p>0,
\]
the formula
\[
    d\nu(x):=\frac{d(x)}{p}\,d\mu(x)
\]
defines a probability measure on $(\Omega,\mathcal A)$.

Let $V(F)=\{1,\dots,n\}$. In the homomorphism-density integral for
$\operatorname{Wh}(F)$, integrate first over the $n$ newly added leaves.
For a fixed value $x_i$ of the original vertex $i$, its private leaf
contributes the factor
\[
    \int_\Omega W(x_i,y)\,d\mu(y)=d(x_i).
\]
Consequently,
\[
\begin{aligned}
    t(\operatorname{Wh}(F),W)
    &=
    \int_{\Omega^n}
       \left(
         \prod_{ij\in E(F)}W(x_i,x_j)
       \right)
       \left(
         \prod_{i=1}^n d(x_i)
       \right)
       d\mu(x_1)\cdots d\mu(x_n)
    \\
    &=
    p^n
    \int_{\Omega^n}
       \prod_{ij\in E(F)}W(x_i,x_j)
       \,d\nu(x_1)\cdots d\nu(x_n).
\end{aligned}
\]
Define
\[
    t_\nu(F,W)
    :=
    \int_{\Omega^{V(F)}}
       \prod_{ij\in E(F)}W(x_i,x_j)
       \prod_{i\in V(F)}d\nu(x_i).
\]
Thus
\begin{equation}\label{eq:whisker-weighted}
    t(\operatorname{Wh}(F),W)
    =
    p^n t_\nu(F,W).
\end{equation}

We next compute the edge density of $W$ with respect to the
probability measure $\nu$. Put
\[
    p_\nu
    :=
    \int_{\Omega^2}W(x,y)\,d\nu(x)d\nu(y).
\]
Then
\[
\begin{aligned}
    p_\nu
    &=
    \frac{1}{p^2}
    \int_{\Omega^2}
       W(x,y)d(x)d(y)\,d\mu(x)d\mu(y)
    \\
    &=
    \frac{t(P_4,W)}{p^2}.
\end{aligned}
\]
By Lemma~\ref{lem:P4},
\[
    t(P_4,W)\ge p^3,
\]
and hence
\begin{equation}\label{eq:pnu-ge-p}
    p_\nu\ge p.
\end{equation}
In particular,
\[
    1\ge p_\nu\ge p\ge p_0,
\]
where $p_\nu\le1$ follows from $0\le W\le1$ and the fact that $\nu$
is a probability measure.

The same measurable kernel $W$ is a graphon on the probability space
$(\Omega,\mathcal A,\nu)$: changing the measure does not alter its
measurability, symmetry, or pointwise range.  Its edge density on this
space is $p_\nu$, and its $F$-density is exactly $t_\nu(F,W)$.
The arbitrary-probability-space hypothesis for $F$ therefore applies
directly and gives
\[
    t_\nu(F,W)
    \ge
    \Phi_F(p_\nu).
\]
Since $\Phi_F$ is nondecreasing on $[p_0,1]$ and
$p_\nu\ge p$, we obtain
\[
    t_\nu(F,W)
    \ge
    \Phi_F(p_\nu)
    \ge
    \Phi_F(p).
\]
Together with \eqref{eq:whisker-weighted}, this yields
\begin{equation}\label{eq:whisker-main}
    t(\operatorname{Wh}(F),W)
    \ge
    p^n\Phi_F(p).
\end{equation}

It remains to identify the right-hand side with the chromatic-polynomial
target for $\operatorname{Wh}(F)$.

Fix a positive integer $k$. Once a proper $k$-colouring of $F$ has
been chosen, each newly added leaf has exactly $k-1$ available colours,
independently of all other leaves. Hence
\[
    \chi_{\operatorname{Wh}(F)}(k)
    =
    \chi_F(k)(k-1)^n.
\]
Since this holds for every positive integer $k$, it is a polynomial
identity:
\begin{equation}\label{eq:whisker-chromatic}
    \chi_{\operatorname{Wh}(F)}(x)
    =
    \chi_F(x)(x-1)^n.
\end{equation}

Put
\[
    q:=1-p.
\]
For $p<1$, using
$v(\operatorname{Wh}(F))=2n$ and
\eqref{eq:whisker-chromatic}, we obtain
\[
\begin{aligned}
    \Phi_{\operatorname{Wh}(F)}(p)
    &=
    q^{2n}
    \chi_{\operatorname{Wh}(F)}(q^{-1})
    \\
    &=
    q^{2n}
    \chi_F(q^{-1})
    (q^{-1}-1)^n
    \\
    &=
    q^{2n}
    \chi_F(q^{-1})
    \left(\frac{p}{q}\right)^n
    \\
    &=
    p^n q^n \chi_F(q^{-1})
    \\
    &=
    p^n\Phi_F(p).
\end{aligned}
\]
The same identity holds at $p=1$ by continuity. Combining this with
\eqref{eq:whisker-main} gives
\[
    t(\operatorname{Wh}(F),W)
    \ge
    \Phi_{\operatorname{Wh}(F)}(p),
\]
as required.
\end{proof}

\begin{corollary}[Preservation of the chromatic threshold]
Let $F$ be a graph with
\[
    r:=\chi(F)\ge3.
\]
Assume that
\[
    t(F,W)\ge\Phi_F(p)
\]
for every graphon $W$ with
\[
    p=t(K_2,W)\ge 1-\frac1{r-1},
\]
and assume that $\Phi_F$ is nondecreasing on
\[
    \left[1-\frac1{r-1},1\right].
\]
Then
\[
    t(\operatorname{Wh}(F),W)
    \ge
    \Phi_{\operatorname{Wh}(F)}(p)
\]
for every graphon $W$ with
\[
    p\ge1-\frac1{r-1}.
\]
Moreover,
\[
    \chi(\operatorname{Wh}(F))=\chi(F)=r.
\]
Thus full whiskering preserves both the Goodman-type inequality and
the required threshold $p\ge1-1/(\chi(F)-1)$.
\end{corollary}

\begin{proof}
The density inequality follows from
Theorem~\ref{thm:whiskering} with
\[
    p_0=1-\frac1{r-1}>0.
\]

Since $F$ is a subgraph of $\operatorname{Wh}(F)$,
\[
    \chi(\operatorname{Wh}(F))\ge\chi(F)=r.
\]
Conversely, take any proper $r$-colouring of $F$. Since $r\ge2$, each
new leaf can be assigned any colour different from the colour of its
unique neighbour. Hence this colouring extends to a proper $r$-colouring
of $\operatorname{Wh}(F)$, so
\[
    \chi(\operatorname{Wh}(F))\le r.
\]
Therefore
\[
    \chi(\operatorname{Wh}(F))=r.
\]
\end{proof}

\begin{corollary}[Atlas 94]\label{cor:atlas94}
NetworkX Atlas graph $94$, with graph6 code \texttt{E@dW}, is positive:
for every graphon $W$ of edge density $p\in[1/2,1]$,
\[
  t(F_{94},W)\geq p^4(2p-1)=\Phi_{F_{94}}(p).
\]
Its order, size, and chromatic number are respectively $6$, $6$, and $3$,
and
\[
  \chi_{F_{94}}(x)=x(x-1)^4(x-2).
\]
\end{corollary}

\begin{proof}
Let $F=K_3$.  The accepted triangle (equivalently, length-three
odd-cycle) bound states that, on every probability space,
\[
  t(K_3,W)\geq p(2p-1)=\Phi_{K_3}(p)
  \qquad (1/2\leq p\leq1).
\]
The polynomial $p(2p-1)$ is nonnegative and strictly increasing on
$[1/2,1]$, since its derivative is $4p-1>0$ there.  Thus
Theorem~\ref{thm:whiskering} applies with $n=3$ and $p_0=1/2$.

In the labelling with triangle edges $01,02,12$ and private leaf edges
$03,14,25$, the result is $\operatorname{Wh}(K_3)$.  Atlas 94 has
NetworkX edge list
\[
  04,\ 15,\ 23,\ 34,\ 35,\ 45.
\]
The relabelling
\[
  0\mapsto4,\quad 1\mapsto3,\quad 2\mapsto5,\quad
  3\mapsto0,\quad 4\mapsto2,\quad 5\mapsto1
\]
carries the first edge list to the second.  Hence Atlas 94 is exactly
the full whiskering of the triangle.  The chromatic identity gives
\[
  \chi_{F_{94}}(x)
  =x(x-1)(x-2)(x-1)^3
  =x(x-1)^4(x-2),
\]
and the density and target identities give
\[
  t(F_{94},W)\geq p^3\Phi_{K_3}(p)=p^4(2p-1)
  =\Phi_{F_{94}}(p).
\]
At $p=1/2$, the target vanishes (and the balanced bipartite graphon
gives equality because the graph contains a triangle).  At $p=1$,
$W=1$ almost everywhere and both sides equal $1$.
\end{proof}

This is exhaustive for the scoped Atlas catalogue on at most six
vertices.  A full whiskering in that range has a base of order at most
three.  If its result is non-bipartite, then the base is non-bipartite;
among graphs of order at most three this forces a triangle component.
Any additional isolated base vertex would become an excluded isolated
edge component.  Thus $K_3$ is the only eligible base and Atlas 94 is
the only scoped full-whiskering row.

\begin{remark}[Iteration]
Suppose, in addition, that $\Phi_F\ge0$ and is nondecreasing on the
relevant interval. Since
\[
    \Phi_{\operatorname{Wh}(F)}(p)=p^{v(F)}\Phi_F(p),
\]
the same two properties hold for $\Phi_{\operatorname{Wh}(F)}$.
Consequently the operation may be iterated.

If
\[
    F_0:=F,
    \qquad
    F_{j+1}:=\operatorname{Wh}(F_j),
\]
then
\[
    v(F_j)=2^jv(F)
\]
and
\[
    \Phi_{F_j}(p)
    =
    p^{(2^j-1)v(F)}\Phi_F(p).
\]
Thus every iterated full whiskering satisfies the same
Goodman-type inequality on the same edge-density interval.
\end{remark}

\subsection{Inputs, hidden-assumption audit, and exact limitations}

The general closure theorem uses only Tonelli--Fubini, the three-edge
path inequality proved above, degree-biasing of the original probability
measure, and the assumed inequality and monotonicity for $F$.
Corollary~\ref{cor:atlas94} additionally uses exactly the accepted
triangle bound.  No regularity, minimum-degree hypothesis, finite-step
representation, or injective-copy interpretation is used.  Points with
$d_W(x)=0$ are handled by an almost-everywhere convention justified
before H\"older is applied, and the proof never divides by $d_W(x)$
pointwise outside that convention.

The hypothesis $p_0>0$ is essential to the degree-biased probability
measure as written.  It is automatic for the required threshold when
$\chi(F)\ge3$.  The endpoint $p=1$ is handled by the polynomial
continuation and directly by the same measure-change proof; no
substitution of $1/(1-p)$ is made there.

Full whiskering means one distinct private leaf at every vertex of
$F$.  The theorem does not cover a leaf at only some vertices, several
leaves at a vertex, shared leaves, or arbitrary pendant-tree attachment.
It also requires monotonicity of $\Phi_F$ on the whole density interval;
positivity of $F$ alone does not imply that monotonicity.  Consequently
this note is not a general leaf-attachment closure theorem.

%%%%%%%%%%%%%%%%%%%%%%%%%%%%%%%%%%%%%%%%%%%%%%%%%%%%%%%%%%%%%%%%%%%%%%%%%%%%%%%
\appendix
\section{What the Lean formalization actually proves}
%%%%%%%%%%%%%%%%%%%%%%%%%%%%%%%%%%%%%%%%%%%%%%%%%%%%%%%%%%%%%%%%%%%%%%%%%%%%%%%

Corollary~\ref{cor:atlas94} is formalized: Atlas $94$ is \texttt{verified},
as \texttt{satisfiesLowerBound\_94} in
\texttt{lean/Taeyoung/Methods/Whisker94.lean}.  One difference of substance.

\subsection*{Theorem~\ref{thm:whiskering} is not stated in general}

The change of measure is formalized in full and in general
(\texttt{Methods/DegreeBias.lean}): $d\nu=(d_W/p)\,d\mu$ is a probability
measure, $W$ reread on $(\Omega,\nu)$ is a graphon, its densities are the
degree-weighted integrals \eqref{eq:whisker-weighted}, and
\eqref{eq:pnu-ge-p} --- the biased edge density is at least $p$, by
Lemma~\ref{lem:P4} --- is \texttt{le\_cliqueDensity\_degreeGraphon}.  That is
the whole analytic content of the theorem.

What is not formalized is the combinatorial half: $\operatorname{Wh}(F)$ is not
constructed as a graph operation, so the identity
$t(\operatorname{Wh}(F),W)=p^{n}\,t(F,W_\nu)$ is not available for a general
$F$.  It is supplied per row, by the peeling identity for that row's own
graph, run on both sides --- six coordinates and three for Atlas $94$.  The
hypotheses of Theorem~\ref{thm:whiskering} are then discharged concretely:
Goodman for $F=K_3$, and monotonicity of $z\mapsto z(2z-1)$ on $[1/2,1]$ in
place of the abstract hypothesis (2).

\end{document}
